\chapter{Explicit Reflectors and the Composition Theorem}
\label{ch10:chapter}

\section{Purpose and status of this chapter}
\label{ch10:sec:purpose}

The preceding chapters established four finite-context facts about the marked
context factor: the transverse partition families generate genuine Grothendieck
topologies $J_A$ and $J_R$ on the sixteen-object lattice $L$; both topologies
are subcanonical; the associated-sheaf modalities $M_A$ and $M_R$ are distinct
and neither is trivial, witnessed by the explicit presheaves $F_A$ and $F_R$;
and the common refinement of the two global covers is the four-singleton cover.
Each of those results concerns one topology at a time, or the raw combinatorics
of the covers. None of them says what happens when the two modalities are
\emph{composed}.

This chapter closes that question for the finite context factor. The answer is
exact and, from the standpoint of the program's earlier informal hopes,
corrective: the two modalities commute up to unit-compatible canonical natural
isomorphism, and either composite is naturally isomorphic to sheafification for
the joined topology $J_A \vee J_R$, which is itself the topology generated by
the four-singleton partition. The proof is fully constructive. Rather than
invoking sheafification as an abstract left adjoint and reasoning about it
indirectly, we build the associated-sheaf reflector for each partition topology
as an explicit product-over-blocks formula, verify the sheaf condition and the
universal property directly against the definitions, identify the join of the
two topologies by a two-sided sieve argument, and then exhibit the composite of
the two explicit reflectors as a flattening of nested products onto atomic
products. Every step is elementary; the value of the chapter is that no step is
skipped.

\begin{boundary}
\label{ch10:bnd:status}
Everything in this chapter transcribes and expands the registered artifact
\texttt{MODALITY\_COMPOSITION\_PROOF\_v0.1.md}, whose status is
\textsc{internal proof candidate / review required}, admission Ring~4,
\textsc{not licensed}, disclosure internal only. Drafting this chapter promotes
nothing: the composition theorem has not been reviewed by an independent
category theorist, has not been formalized in a proof assistant, and licenses
no empirical, probabilistic, product, or safety claim. The exhaustive
executable corroboration described in Chapter~11 is corroboration, not proof.
The result is scoped to $\PSh(L, \Set)$; no statement in this chapter lifts it
to $\PSh(\cC \times L)$, includes the observational topology on $\cC$, or
interacts with the legacy comparison of Chapter~12.
\end{boundary}

\section{Standing data}
\label{ch10:sec:standing}

We fix the data used throughout, in the exact form assumed by the registered
proof.

\begin{notation}
\label{ch10:not:standing}
$\Omega = \{(0,0), (0,1), (1,0), (1,1)\}$ is the four-atom set; its elements
are written $\omega_{ij} = (i,j)$. The context lattice is
$L = \mathcal{P}(\Omega)$, regarded as a category with one arrow $U \to V$
exactly when $U \subseteq V$. A presheaf is a functor
$F : L^{\mathrm{op}} \to \Set$, small in a fixed universe. The two transverse
partitions of $\Omega$ are
\[
A_i = \{(i,0), (i,1)\}, \qquad R_j = \{(0,j), (1,j)\}, \qquad i, j \in \{0,1\},
\]
and $Q$ is the partition of $\Omega$ into its four singletons. For a partition
$P \in \{A, R, Q\}$ and a nonempty context $V$, the generating family is
\[
K_P(V) = \{\, V \cap B \to V \mid B \text{ a block of } P \,\},
\]
retaining an empty piece $\varnothing \to V$ whenever an intersection is empty.
At the empty context the generator is the null (empty) family, which covers.
$J_A$, $J_R$, $J_{\mathrm{atom}}$ are the topologies these bases generate; the
publication-layer name for $J_{\mathrm{atom}}$ is $J_{\mathrm{singleton}}$,
since the label does not assert the standard category-theoretic atomic
topology.
\end{notation}

The registered proof assumes exactly five things: ordinary small $\Set$-valued
presheaves in a fixed universe; the powerset inclusion category $L$; the stated
generating families with the null-cover convention; the ordinary
Grothendieck-topology and associated-sheaf definitions
\cite{maclane, stacks00ZC}; and the transversality identity
\begin{equation}
\label{ch10:eq:transverse}
A_i \cap R_j = \{(i,j)\},
\end{equation}
which is immediate from the definitions of the blocks. No legacy theorem, no
semantic interpretation of the ACBE or RPM marks, and no empirical premise is
used anywhere below. Pullbacks in $L$ are intersections: the pullback of
$U \to V$ along $W \to V$ is $U \cap W$.

Two elementary facts about sieves are used repeatedly and are recorded here
once.

\begin{lemma}
\label{ch10:lem:supersieve}
Let $J$ be a Grothendieck topology on a category, $S \in J(C)$ a covering
sieve, and $T \supseteq S$ any sieve on $C$. Then $T \in J(C)$.
\end{lemma}

\begin{proof}
Apply the local-character axiom (J3) to $T$ with respect to the covering sieve
$S$: for each $f : D \to C$ in $S$ we have $f \in T$, hence
$f^{*}T \ni \id_D$, so $f^{*}T$ is the maximal sieve on $D$ and is covering by
(J1). Therefore $T \in J(C)$.
\end{proof}

\begin{lemma}
\label{ch10:lem:coarser}
If $J \le J'$ (every $J$-covering sieve is $J'$-covering), then every
$J'$-sheaf is a $J$-sheaf.
\end{lemma}

\begin{proof}
The sheaf condition for a presheaf quantifies over covering sieves; a
$J'$-sheaf satisfies the gluing bijection for every sieve in $J'$, in
particular for every sieve in $J \subseteq J'$.
\end{proof}

Finally, we recall the basis criterion used throughout: the generating families
above form a pullback-stable coverage basis (pullback along $U \subseteq V$
carries $K_P(V)$ to $K_P(U)$ as an indexed family, and to identity arrows on
the empty object when $U = \varnothing$; refinement of a member by the same
partition regenerates the same sieve because distinct blocks are disjoint), and
a presheaf is a sheaf for the generated topology if and only if it satisfies
the gluing condition for the basis families \cite{stacks00ZC}. This was proved
in full when the topologies were first constructed; we use it without further
comment.

\section{The one-partition reflector}
\label{ch10:sec:reflector}

The composition theorem requires knowing sheafification \emph{exactly}, not
just abstractly as a left adjoint. On the finite lattice $L$ with a partition
coverage, sheafification admits a closed form: a product over the blocks the
context meets.

\begin{construction}[One-partition reflector]
\label{ch10:con:reflector}
Fix a partition $P \in \{A, R, Q\}$. For a context $V$, let
\[
I_P(V) = \{\, B \in P \mid V \cap B \neq \varnothing \,\}
\]
be the set of blocks that $V$ meets. For every presheaf $F$, define a presheaf
$R_P F$ by
\begin{equation}
\label{ch10:eq:reflector}
(R_P F)(V) \;=\; \prod_{B \in I_P(V)} F(V \cap B),
\end{equation}
where the empty product is the singleton $\mathbf{1}$. In particular
\begin{equation}
\label{ch10:eq:emptyclause}
(R_P F)(\varnothing) = \mathbf{1}
\end{equation}
even when $F(\varnothing)$ has zero, two, or more elements. For an inclusion
$U \subseteq V$, the restriction map
$(R_P F)(V) \to (R_P F)(U)$ \emph{drops} the coordinates $B$ with
$U \cap B = \varnothing$ and applies the restriction map
$F(V \cap B) \to F(U \cap B)$ of $F$ on each retained coordinate. The unit is
the natural transformation
\begin{equation}
\label{ch10:eq:unit}
\eta_{P, F, V} : F(V) \to (R_P F)(V), \qquad
x \mapsto \bigl( x|_{V \cap B} \bigr)_{B \in I_P(V)},
\end{equation}
where $x|_{V \cap B}$ denotes the image of $x$ under the restriction map of
$F$ along $V \cap B \subseteq V$.
\end{construction}

\begin{proposition}[Well-definedness]
\label{ch10:prop:welldef}
$R_P F$ is a presheaf, $R_P$ is a functor on $\PSh(L, \Set)$, and $\eta_P$ is a
natural transformation $\id \Rightarrow R_P$.
\end{proposition}

\begin{proof}
\emph{Functor laws for $R_P F$.} Restriction along $\id_V$ retains every
coordinate of $I_P(V)$ and applies $F$-restriction along identities, which is
the identity by functoriality of $F$. For $W \subseteq U \subseteq V$: a
coordinate $B$ survives the composite $V \to U \to W$ restriction exactly when
$W \cap B \neq \varnothing$, which (since $W \subseteq U$ forces
$U \cap B \neq \varnothing$ as well) is the same survival condition as for the
one-step restriction $V \to W$; on each surviving coordinate the two routes
apply $F(V \cap B) \to F(U \cap B) \to F(W \cap B)$ and
$F(V \cap B) \to F(W \cap B)$, equal by functoriality of $F$. Coordinates
dropped at the first stage stay dropped, coordinatewise agreement gives
equality of the maps.

\emph{Functoriality of $R_P$.} A presheaf map $\tau : F \to F'$ induces
$(R_P \tau)_V = \prod_{B \in I_P(V)} \tau_{V \cap B}$; this commutes with the
restriction maps of $R_P F$ and $R_P F'$ because $\tau$ is natural at each
$V \cap B \supseteq U \cap B$, and it preserves identities and composites
coordinatewise.

\emph{Naturality of the unit.} For $U \subseteq V$ and $x \in F(V)$, both
routes around the naturality square send $x$ to the tuple
$\bigl( x|_{U \cap B} \bigr)_{B \in I_P(U)}$: going right then down restricts
each coordinate $x|_{V \cap B}$ further to $U \cap B$ and drops the
coordinates $U$ misses; going down then right restricts $x$ to $U$ first and
then forms the tuple. Both computations are composites of restriction maps of
$F$ along the same inclusions, hence equal. Naturality of $\eta_P$ in $F$
(that is, $\eta_{P, F'} \circ \tau = (R_P \tau) \circ \eta_{P, F}$ for
$\tau : F \to F'$) is the componentwise naturality of $\tau$ along the
inclusions $V \cap B \subseteq V$.
\end{proof}

\begin{remark}[Two structural observations]
\label{ch10:rem:structure}
Two features of formula \eqref{ch10:eq:reflector} do all the later work.
First, the value at a context that meets exactly one block is canonically the
value of $F$ at the intersection: if $V \neq \varnothing$ and
$V \subseteq B_0$ for a block $B_0$, then $I_P(V) = \{B_0\}$ and
$(R_P F)(V) = F(V \cap B_0) = F(V)$ via the one-factor product, and under this
identification the unit $\eta_{P, F, V}$ is the identity map of $F(V)$. We
call this the \emph{singleton-factor bijection}; it is the engine of the
uniqueness argument in Lemma~\ref{ch10:lem:universal}. Second, the value at
the empty context is forcibly the singleton, clause
\eqref{ch10:eq:emptyclause}, regardless of $F(\varnothing)$. This is not an
arbitrary convention: the null cover of the empty context has exactly one
matching family (the empty one), so every sheaf for any of the three
topologies must have a one-element value at $\varnothing$
(Lemma~\ref{ch10:lem:emptyterminal} below), and the reflector must therefore
correct a nonterminal empty-context value. The executable fixture of
Chapter~11 that exercises exactly this clause is a constant two-section
presheaf whose empty-context value is deliberately nonterminal.
\end{remark}

\begin{lemma}
\label{ch10:lem:emptyterminal}
Let $J \in \{J_A, J_R, J_{\mathrm{atom}}\}$. Every $J$-sheaf $G$ has
$G(\varnothing)$ a one-element set.
\end{lemma}

\begin{proof}
The null family covers $\varnothing$ in each basis, hence the sieve it
generates (the empty sieve) is covering. A matching family for the empty
family is a family of sections indexed by the empty set; there is exactly one
such family. The sheaf condition demands that the canonical map from
$G(\varnothing)$ to this one-element set of matching families be a bijection,
so $G(\varnothing)$ is a singleton.
\end{proof}

\begin{example}[The reflector in action]
\label{ch10:ex:action}
Three computations calibrate the construction; the appendices tabulate the
witness values used here.
\begin{enumerate}
\item \emph{Representables are fixed.} For $yW$ with $W \in L$, the
subcanonicity theorem (Theorem NT2-2, proved earlier in the monograph) says
$yW$ is already a $J_P$-sheaf; by Corollary~\ref{ch10:cor:fixed} below, the
unit $yW \to R_P\, yW$ is an isomorphism. Concretely: a matching tuple of
compatible sections over the pieces of $V$ exists exactly when every piece
lies in $W$, which by joint exhaustiveness of the partition happens exactly
when $V \subseteq W$.
\item \emph{The ACBE reflector annihilates the RPM witness.} The RPM witness
$F_R$ assigns to $V$ the bit assignments on the complete R-blocks contained
in $V$. For every context $V$ and every $i$, the piece $V \cap A_i$ is a
subset of $A_i$ and therefore contains no complete R-block, since
$R_j \subseteq A_i$ is impossible by transversality
\eqref{ch10:eq:transverse}. Hence $F_R(V \cap A_i) = \mathbf{1}$ for every
piece, and $(R_A F_R)(V) = \mathbf{1}$ at every context: $M_A F_R$ is the
terminal presheaf. Symmetrically $M_R F_A$ is terminal. The modalities are not
merely non-identity; each one can destroy all the information carried by the
other's witness.
\item \emph{The empty-context correction.} Let $C$ be the constant presheaf
with $C(V) = \{0, 1\}$ for all $V$, all restrictions the identity. Then
$(R_P C)(\varnothing) = \mathbf{1} \neq C(\varnothing)$: the unit is not
injective at the empty context, and $C$ is not a $J_P$-sheaf precisely because
it fails the null-cover condition. At nonempty $V$, $(R_P C)(V)$ is a product
of two-element sets, one per block $V$ meets; for $V = \Omega$ and $P = A$
this is a four-element set, so the reflector can also \emph{enlarge} values.
\end{enumerate}
\end{example}

\begin{proposition}[Limit preservation]
\label{ch10:prop:lex}
$R_P$ preserves all small limits of presheaves; in particular it is left
exact.
\end{proposition}

\begin{proof}
Limits in $\PSh(L, \Set)$ are computed pointwise. At each context $V$, formula
\eqref{ch10:eq:reflector} exhibits $(R_P F)(V)$ as a finite product of the
evaluation functors $F \mapsto F(V \cap B)$, each of which preserves limits
(evaluation is limit-preserving pointwise), and finite products of
limit-preserving functors preserve limits. Left exactness is the finite-limit
case. This is consistent with, and here independent of, the general fact that
sheafification is left exact \cite{stacks00W1}; we will only conclude that
$R_P$ \emph{is} sheafification in Lemma~\ref{ch10:lem:universal}.
\end{proof}

\section{The sheaf lemma}
\label{ch10:sec:sheaf}

\begin{lemma}[$R_P F$ is a $J_P$-sheaf]
\label{ch10:lem:sheafness}
For every presheaf $F$ and every $P \in \{A, R, Q\}$, the presheaf $R_P F$ is
a sheaf for $J_P$.
\end{lemma}

\begin{proof}
Because the generating families form a coverage basis, it suffices to verify
the gluing bijection for the basis families \cite{stacks00ZC}. There are two
cases.

\emph{The empty context.} The null cover requires exactly one matching family
and exactly one section. The empty family has exactly one matching family (the
empty one), and $(R_P F)(\varnothing) = \mathbf{1}$ by clause
\eqref{ch10:eq:emptyclause}, so the canonical map is a bijection between
one-element sets.

\emph{A nonempty context $V$ with cover $K_P(V)$.} We first compute the value
of $R_P F$ on each piece, then dispose of the compatibility constraints, then
exhibit the gluing bijection.

\emph{Values on pieces.} If $V \cap B \neq \varnothing$, then
$I_P(V \cap B) = \{B\}$: the piece is contained in the single block $B$, and
blocks are pairwise disjoint, so it meets no other block. Hence
\[
(R_P F)(V \cap B) = F\bigl((V \cap B) \cap B\bigr) = F(V \cap B),
\]
the singleton-factor bijection of Remark~\ref{ch10:rem:structure}. If
$V \cap B = \varnothing$, the piece is the empty context, where
$(R_P F)(\varnothing) = \mathbf{1}$.

\emph{The disjoint-overlap argument.} Let $B \neq B'$ be distinct blocks. The
overlap of the corresponding pieces is the pullback
\[
(V \cap B) \cap (V \cap B') = V \cap (B \cap B') = V \cap \varnothing
= \varnothing,
\]
since distinct blocks of a partition are disjoint. Compatibility of a matching
family on this overlap is the requirement that the two restricted sections
agree in $(R_P F)(\varnothing) = \mathbf{1}$ --- a one-element set, in which
any two elements are equal. Compatibility is therefore \emph{automatic} for
every pair of distinct pieces. More generally, any arrow $U \to V$ that
factors through two distinct pieces satisfies
$U \subseteq (V \cap B) \cap (V \cap B') = \varnothing$, so all overlap
conditions in the sieve generated by $K_P(V)$ reduce to conditions in the
one-element set. Compatibility of a piece with itself is trivially the
requirement that a section agree with itself. Empty pieces impose no
constraint at all: a matching family assigns to the arrow
$\varnothing \to V$ the unique element of $\mathbf{1}$.

\emph{The gluing bijection.} By the two preceding paragraphs, a matching
family for $K_P(V)$ is precisely a free choice of one element of
$F(V \cap B)$ for each block $B \in I_P(V)$ (the nonempty pieces), together
with the forced unique elements on empty pieces. The set of matching families
is therefore canonically
\[
\prod_{B \in I_P(V)} F(V \cap B) \;=\; (R_P F)(V).
\]
The canonical gluing map sends a section $s \in (R_P F)(V)$ to the family of
its restrictions to the pieces. On the nonempty piece $V \cap B$, the
restriction map of $R_P F$ retains exactly the coordinate $B$ (all other
blocks meet $V \cap B$ in $\varnothing$) and applies the identity
$F(V \cap B) \to F(V \cap B)$; that is, restriction to the piece $V \cap B$
is projection onto the $B$-coordinate. The gluing map is therefore the
canonical reindexing of the product by its own coordinates, which is a
bijection.
\end{proof}

\begin{remark}
\label{ch10:rem:whyempty}
The proof shows why clause \eqref{ch10:eq:emptyclause} is load-bearing twice
over. If one had instead set $(R_P F)(\varnothing) = F(\varnothing)$, the
null-cover case would fail whenever $F(\varnothing)$ is not a singleton, and
the disjoint-overlap argument would fail as well: compatibility on empty
overlaps would become a genuine constraint (agreement in $F(\varnothing)$),
and the matching-family set would no longer be the full product. Both failure
modes are exercised by the constant two-section fixture of
Example~\ref{ch10:ex:action}(3), and both are checked executably in
Chapter~11.
\end{remark}

\section{The universal property}
\label{ch10:sec:universal}

Lemma~\ref{ch10:lem:sheafness} shows $R_P F$ is \emph{a} sheaf receiving a map
from $F$. The next lemma shows it is \emph{the} associated sheaf: the unit is
universal among maps from $F$ to $J_P$-sheaves. This is the step that
identifies the explicit formula with the abstract modality, and it is proved
directly rather than by appeal to a general sheafification construction.

\begin{lemma}[Universal property]
\label{ch10:lem:universal}
For every presheaf $F$ and every $J_P$-sheaf $G$, composition with the unit
\eqref{ch10:eq:unit} is a bijection
\[
\Hom_{\Sh(L, J_P)}(R_P F, G) \;\xrightarrow{\;\cong\;}\; \Hom_{\PSh(L)}(F, G),
\qquad \beta \mapsto \beta \circ \eta_{P, F},
\]
natural in $F$ and $G$. Consequently $\eta_P : \id \Rightarrow R_P$ exhibits
$R_P$ as the associated-sheaf reflector for $J_P$, left adjoint to the
inclusion $\Sh(L, J_P) \hookrightarrow \PSh(L, \Set)$.
\end{lemma}

\begin{proof}
Fix $\alpha : F \to G$ with $G$ a $J_P$-sheaf. We construct a factorization,
prove it is a presheaf map, prove it factors $\alpha$, prove it is unique, and
then address naturality.

\emph{Existence.} Let $V$ be a context and
$(x_B)_{B \in I_P(V)} \in (R_P F)(V)$ a tuple. Apply $\alpha$ componentwise:
for each $B \in I_P(V)$, the element
$\alpha_{V \cap B}(x_B) \in G(V \cap B)$ is a section of $G$ over the piece
$V \cap B$; assign to each empty piece the unique element of
$G(\varnothing) = \mathbf{1}$ (Lemma~\ref{ch10:lem:emptyterminal}). These
local sections form a matching family for $K_P(V)$: every overlap of distinct
pieces is the empty context by the disjoint-overlap computation of
Lemma~\ref{ch10:lem:sheafness}, and $G(\varnothing)$ is a one-element set, so
all compatibility conditions hold automatically. Since $G$ is a $J_P$-sheaf,
this matching family glues to a \emph{unique} section of $G(V)$; define
\[
\bar{\alpha}_V\bigl((x_B)_B\bigr) \in G(V)
\]
to be that glued section, characterized by
$\bar{\alpha}_V((x_B))|_{V \cap B} = \alpha_{V \cap B}(x_B)$ for every
$B \in I_P(V)$. At $V = \varnothing$, define $\bar{\alpha}_\varnothing$ to be
the unique map $\mathbf{1} \to G(\varnothing)$ between one-element sets.

\emph{Naturality of $\bar{\alpha}$ in $V$.} Let $U \subseteq V$ and
$(x_B) \in (R_P F)(V)$; write $s = \bar{\alpha}_V((x_B))$. We must show
$s|_U = \bar{\alpha}_U\bigl((x_B)|_U\bigr)$, where $(x_B)|_U$ is the
restricted tuple $\bigl( x_B|_{U \cap B} \bigr)_{B \in I_P(U)}$. Both sides
are sections of $G(U)$; by the uniqueness half of the sheaf condition for the
cover $K_P(U)$ of $U$ (separatedness), it suffices to show they agree on each
piece $U \cap B$ with $B \in I_P(U)$. The right side restricts to
$\alpha_{U \cap B}(x_B|_{U \cap B})$ by its defining property. The left side
restricts to
\[
(s|_U)|_{U \cap B} = s|_{U \cap B} = \bigl( s|_{V \cap B} \bigr)|_{U \cap B}
= \bigl( \alpha_{V \cap B}(x_B) \bigr)|_{U \cap B}
= \alpha_{U \cap B}\bigl( x_B|_{U \cap B} \bigr),
\]
using functoriality of restriction in $G$ along
$U \cap B \subseteq V \cap B \subseteq V$, the defining property of $s$, and
naturality of $\alpha$ along $U \cap B \subseteq V \cap B$. The two sections
agree on every piece of a cover of $U$, hence are equal. (When
$U = \varnothing$ both sides lie in the one-element set $G(\varnothing)$.)

\emph{Factorization.} For $x \in F(V)$,
$\bar{\alpha}_V(\eta_{P, F, V}(x)) = \bar{\alpha}_V\bigl((x|_{V \cap B})_B\bigr)$
is the unique section of $G(V)$ restricting on each piece to
$\alpha_{V \cap B}(x|_{V \cap B})$. By naturality of $\alpha$, the section
$\alpha_V(x)$ has exactly these restrictions:
$\alpha_V(x)|_{V \cap B} = \alpha_{V \cap B}(x|_{V \cap B})$. Separatedness of
$G$ for $K_P(V)$ forces $\bar{\alpha}_V(\eta_{P, F, V}(x)) = \alpha_V(x)$,
i.e. $\bar{\alpha} \circ \eta_{P, F} = \alpha$.

\emph{Uniqueness via the singleton-factor bijection.} Let
$\beta : R_P F \to G$ be any presheaf map with
$\beta \circ \eta_{P, F} = \alpha$. First consider a context $W$ that is a
nonempty subset of a single block $B$. There the unit
$\eta_{P, F, W} : F(W) \to (R_P F)(W)$ is the singleton-factor bijection of
Remark~\ref{ch10:rem:structure} (one coordinate, identity map). Hence on such
contexts $\beta$ is completely determined:
$\beta_W = \alpha_W \circ \eta_{P, F, W}^{-1}$. Now let $V$ be arbitrary
nonempty and $(x_B) \in (R_P F)(V)$. For each $B \in I_P(V)$, naturality of
$\beta$ along $V \cap B \subseteq V$ gives
\[
\beta_V\bigl((x_B)\bigr)\big|_{V \cap B}
= \beta_{V \cap B}\bigl( (x_B)|_{V \cap B} \bigr)
= \beta_{V \cap B}\bigl( \eta_{P, F, V \cap B}(x_B) \bigr)
= \alpha_{V \cap B}(x_B),
\]
where the middle equality uses that restriction of $R_P F$ to the piece
$V \cap B$ is projection onto the $B$-coordinate, which under the
singleton-factor bijection is exactly the unit applied to $x_B$. So
$\beta_V((x_B))$ is a section of $G(V)$ whose restriction to every piece of
$K_P(V)$ agrees with that of $\bar{\alpha}_V((x_B))$; separatedness forces
$\beta_V = \bar{\alpha}_V$. At the empty context both maps land in the
one-element set $G(\varnothing)$ and are equal. Hence $\beta = \bar{\alpha}$,
and the assignment $\alpha \mapsto \bar{\alpha}$ is a two-sided inverse to
composition with the unit.

\emph{Naturality of the bijection.} In $G$: for a sheaf map $g : G \to G'$,
we must check $\overline{g \circ \alpha} = g \circ \bar{\alpha}$; both are
factorizations of $g \circ \alpha$ through the unit, and uniqueness (already
proved, applied with target $G'$) forces equality. In $F$: for a presheaf map
$\tau : F' \to F$, we must check
$\overline{\alpha \circ \tau} = \bar{\alpha} \circ R_P \tau$; the right side
is a sheaf map $R_P F' \to G$ satisfying
$(\bar{\alpha} \circ R_P \tau) \circ \eta_{P, F'}
= \bar{\alpha} \circ \eta_{P, F} \circ \tau = \alpha \circ \tau$
(using naturality of the unit in the presheaf argument,
Proposition~\ref{ch10:prop:welldef}), and uniqueness applied to
$\alpha \circ \tau$ forces equality. Thus the bijection is natural in both
variables, and $R_P$ is left adjoint to the inclusion with unit $\eta_P$.
\end{proof}

\begin{remark}[On the method, with the Stacks citation]
\label{ch10:rem:stacks}
The proof deliberately avoids assuming that a single application of the plus
construction sheafifies an arbitrary presheaf, or that sheafification has any
particular closed form; it verifies the universal property from scratch
against the definitions. The general background --- existence of the
associated sheaf, its universal property, and its exactness properties --- is
recorded in the Stacks Project, Section 7.10, especially Theorem 7.10.10 and
Proposition 7.10.12 \cite{stacks00W1}; the textbook treatment is
\cite[Chapter~III]{maclane}. The point of the direct proof is auditability:
every quantifier in Lemma~\ref{ch10:lem:universal} ranges over finite data
that the executable layer of Chapter~11 can and does enumerate.
\end{remark}

\begin{corollary}[$R_P$ is THE reflector]
\label{ch10:cor:thereflector}
Any left adjoint to the inclusion $\Sh(L, J_P) \hookrightarrow \PSh(L, \Set)$
is naturally isomorphic to $R_P$, compatibly with units. Writing $M_P$ for the
composite of any chosen sheafification with the inclusion, there are canonical
natural isomorphisms $M_A \cong R_A$, $M_R \cong R_R$,
$M_{\mathrm{atom}} \cong R_Q$.
\end{corollary}

\begin{proof}
Uniqueness of left adjoints up to canonical natural isomorphism compatible
with units \cite{maclane}, applied to the adjunction established in
Lemma~\ref{ch10:lem:universal}.
\end{proof}

\begin{corollary}[Fixed objects]
\label{ch10:cor:fixed}
The unit $\eta_{P, F} : F \to R_P F$ is an isomorphism if and only if $F$ is a
$J_P$-sheaf. In particular $R_P$ is idempotent up to canonical isomorphism:
$\eta_{P, R_P F}$ and $R_P(\eta_{P, F})$ are isomorphisms
$R_P F \to R_P R_P F$.
\end{corollary}

\begin{proof}
If $\eta_{P, F}$ is an isomorphism, $F$ is isomorphic to the sheaf $R_P F$
and is a sheaf. Conversely, if $F$ is a $J_P$-sheaf, then $\id_F : F \to F$
factors through the unit by Lemma~\ref{ch10:lem:universal}, giving
$r : R_P F \to F$ with $r \circ \eta = \id$; and both
$\eta \circ r$ and $\id_{R_P F}$ factor $\eta$ through itself, so uniqueness
gives $\eta \circ r = \id$. Idempotency follows since $R_P F$ is a sheaf by
Lemma~\ref{ch10:lem:sheafness}.
\end{proof}

\begin{remark}[Modalities and local operators]
\label{ch10:rem:johnstone}
In topos-theoretic language, each $M_P$ is a left-exact idempotent monad on
the presheaf topos --- what this monograph calls a \emph{modality} --- and
such monads correspond to local operators (Lawvere--Tierney topologies) on
$\PSh(L, \Set)$; for a presheaf topos over a small site these in turn
correspond to Grothendieck topologies on the base. The general theory is in
\cite{johnstone} and \cite[Chapter~V]{maclane}. We use this correspondence
only as orientation; every claim below is proved at the level of the explicit
formulas.
\end{remark}

\section{The join of the two topologies}
\label{ch10:sec:join}

The composite of the two reflectors will turn out to be sheafification for a
single topology: the join. We first identify the join exactly.

\begin{definition}
\label{ch10:def:join}
Order Grothendieck topologies on $L$ by $J \le J'$ iff every $J$-covering
sieve is $J'$-covering. The collection of topologies containing two given
topologies $J_A, J_R$ is nonempty (the maximal topology contains both) and
closed under intersection, so it has a least element: the \emph{join}
$J_A \vee J_R$, the smallest topology containing every $J_A$-cover and every
$J_R$-cover.
\end{definition}

\begin{lemma}[Join identification]
\label{ch10:lem:join}
$J_A \vee J_R = J_{\mathrm{atom}}$, where $J_{\mathrm{atom}}$ is the topology
generated by the singleton partition $Q$ with the null cover at the empty
context.
\end{lemma}

\begin{proof}
We prove the two inequalities separately.

\emph{Step 1: $J_A \vee J_R \le J_{\mathrm{atom}}$ (singletons lie inside
pieces).} It suffices to show that $J_{\mathrm{atom}}$ is a topology
containing $J_A$ and $J_R$; the join, being least among such, is then
contained in it. Fix a nonempty $V$ and consider the sieve $S_A$ generated by
$K_A(V)$. Every singleton $\{\omega\} \subseteq V$ with
$\omega = (i, j)$ satisfies $\{\omega\} \subseteq V \cap A_i$: the singleton
lies inside exactly one ACBE piece. Hence the arrow $\{\omega\} \to V$
factors through a member of $K_A(V)$ and belongs to $S_A$. The sieve $S_Q$
generated by $K_Q(V)$ consists precisely of the arrows factoring through some
singleton piece --- that is, all $U \subseteq V$ with $|U| \le 1$ --- and we
have just shown every generator of $S_Q$ lies in $S_A$, so
$S_Q \subseteq S_A$. Since $S_Q$ is $J_{\mathrm{atom}}$-covering,
Lemma~\ref{ch10:lem:supersieve} makes $S_A$ covering in $J_{\mathrm{atom}}$
as well. The same argument with $R_j$ in place of $A_i$ handles $K_R(V)$. At
the empty context all three bases prescribe the null cover, so the covering
sieves agree there. Every generating cover of $J_A$ and of $J_R$ is therefore
$J_{\mathrm{atom}}$-covering; since $J_A$ (resp.\ $J_R$) is the least
topology containing its generating sieves, $J_A \le J_{\mathrm{atom}}$ and
$J_R \le J_{\mathrm{atom}}$, hence $J_A \vee J_R \le J_{\mathrm{atom}}$.

\emph{Step 2: $J_{\mathrm{atom}} \le J_A \vee J_R$ (composite ACBE-then-RPM
refinement via transversality).} Let $J'$ be any topology containing $J_A$
and $J_R$; we show every $J_{\mathrm{atom}}$-generating sieve is
$J'$-covering, which gives $J_{\mathrm{atom}} \le J'$ and in particular
$J_{\mathrm{atom}} \le J_A \vee J_R$. Fix nonempty $V$ and let
$T = S_Q$ be the sieve on $V$ generated by the singleton family; as noted, $T$
is the set of subsets of $V$ of size at most one. We verify $T \in J'(V)$ by
the local-character axiom (J3) applied along the $J'$-covering sieve $S_A$
generated by $K_A(V)$: it suffices that $g^{*}T \in J'(W)$ for every arrow
$g : W \to V$ in $S_A$.

So let $g : W \to V$ lie in $S_A$; then $W \subseteq V \cap A_i$ for some
$i$. If $W = \varnothing$, then $\varnothing \to V$ lies in $T$ (the empty
set has size zero), so $g^{*}T$ contains $\id_\varnothing$ and is maximal,
hence covering. If $W \neq \varnothing$, compute the pullback sieve:
$g^{*}T = \{\, U \subseteq W \mid U \in T \,\} = \{\, U \subseteq W \mid
|U| \le 1 \,\}$. Now consider the RPM cover $K_R(W)$ of $W$, which is
$J_R$-covering hence $J'$-covering. Its pieces are $W \cap R_j$; since
$W \subseteq A_i$, transversality \eqref{ch10:eq:transverse} gives
\[
W \cap R_j \;\subseteq\; A_i \cap R_j \;=\; \{(i, j)\},
\]
so every RPM piece of $W$ is empty or a singleton --- this is the composite
ACBE-then-RPM refinement collapsing to singletons. The sieve generated by
$K_R(W)$ is the set of subsets of $W$ contained in some piece, i.e. exactly
the subsets of size at most one: every singleton $\{(i, j)\} \subseteq W$ is
contained in the piece $W \cap R_j$, and $\varnothing$ is contained in every
piece. Hence the sieve generated by $K_R(W)$ \emph{equals} $g^{*}T$, which is
therefore $J'$-covering. (Null-cover handling: if some ACBE piece
$V \cap A_i$ is empty, the corresponding arrows contribute the
$W = \varnothing$ case already treated; and at $V = \varnothing$ itself the
$J_{\mathrm{atom}}$ generator is the null family, covering in $J'$ because
$J_A \le J'$ prescribes the same null cover.)

By (J3), $T \in J'(V)$. Every generating sieve of $J_{\mathrm{atom}}$ is thus
$J'$-covering, and since $J_{\mathrm{atom}}$ is the least topology containing
its generators, $J_{\mathrm{atom}} \le J'$. Taking $J' = J_A \vee J_R$
completes Step 2, and the two steps give equality.
\end{proof}

\begin{corollary}
\label{ch10:cor:union}
A sieve on $V$ is $(J_A \vee J_R)$-covering if and only if it contains every
singleton of $V$; equivalently, the joined context coverage is the ordinary
union coverage on the finite powerset lattice $\mathcal{P}(\Omega)$.
\end{corollary}

\begin{proof}
By Lemma~\ref{ch10:lem:join} the join is generated by the singleton families.
A sieve containing every singleton of $V$ contains the generated sieve $S_Q$
and is covering by Lemma~\ref{ch10:lem:supersieve}. Conversely, the
collection of sieves containing all singletons of their target is itself a
Grothendieck topology (the verification is the standard one for the union
coverage on a powerset: maximality is clear; pullback stability holds because
singletons of $U \subseteq V$ are singletons of $V$; local character holds
because a singleton missing from $T$ would pull back along its own inclusion
in the covering sieve to a non-covering sieve on a one-point context, where
covering sieves must contain the identity), and it contains the generators,
hence contains $J_{\mathrm{atom}}$.
\end{proof}

\begin{boundary}
\label{ch10:bnd:atomname}
The internal label $J_{\mathrm{atom}}$ does not assert the standard
category-theoretic \emph{atomic topology} (the topology in which every
nonempty sieve covers); $J_{\mathrm{singleton}}$ is the preferred
publication-layer name. Chapter~11 records the exact-topology enumeration
that corroborates Lemma~\ref{ch10:lem:join}: all $298$ sieves over the
sixteen contexts are enumerated, both presentations (join and singleton) are
saturated to a fixed point, the Grothendieck axioms are checked independently
on both computed topologies, and both report the same $168$ covering sieves.
\end{boundary}

\section{The composition theorem}
\label{ch10:sec:composition}

We can now compute the composite of the two explicit reflectors exactly. The
composite $M_A M_R$ of the modalities means: sheafify for $J_R$, include back
into presheaves, sheafify for $J_A$, include back. By
Corollary~\ref{ch10:cor:thereflector} this composite is canonically naturally
isomorphic to $R_A R_R$, so it suffices to compute the latter.

\begin{theorem}[Modality composition]
\label{ch10:thm:composition}
On $\PSh(L, \Set)$ there are unit-compatible canonical natural isomorphisms
\[
M_A M_R \;\cong\; M_{\mathrm{atom}} \;\cong\; M_R M_A,
\]
and, with Lemma~\ref{ch10:lem:join}, either composite is naturally isomorphic
to sheafification for the joined topology $J_A \vee J_R$. These are natural
isomorphisms, not strict equalities.
\end{theorem}

\begin{proof}
\emph{The objectwise computation.} Fix a presheaf $F$ and a context $V$.
Expand formula \eqref{ch10:eq:reflector} twice:
\[
(R_A R_R F)(V)
= \prod_{B \in I_A(V)} (R_R F)(V \cap B)
= \prod_{\substack{i \\ V \cap A_i \neq \varnothing}}
  \;\prod_{\substack{j \\ V \cap A_i \cap R_j \neq \varnothing}}
  F(V \cap A_i \cap R_j).
\]
By transversality \eqref{ch10:eq:transverse},
\[
V \cap A_i \cap R_j = V \cap \{(i, j)\}
= \begin{cases} \{(i, j)\} & \text{if } (i, j) \in V, \\
\varnothing & \text{otherwise,} \end{cases}
\]
so the inner index condition $V \cap A_i \cap R_j \neq \varnothing$ says
exactly $(i, j) \in V$, and the outer condition
$V \cap A_i \neq \varnothing$ says exactly that some $j$ has
$(i, j) \in V$ --- which is implied by the inner condition and therefore
imposes no further restriction on the pairs that occur. The set of index
pairs $(i, j)$ appearing in the nested product is thus in canonical bijection
with the set of atoms $\omega \in V$, and each factor is $F(\{\omega\})$.

\emph{The flattening and regrouping bijections.} Define
\[
\varphi_{F, V} : (R_A R_R F)(V) \longrightarrow
\prod_{\omega \in V} F(\{\omega\}) = (R_Q F)(V),
\qquad
\varphi_{F, V}\Bigl( \bigl( (x_{i j})_j \bigr)_i \Bigr)
= \bigl( x_{i j} \bigr)_{(i, j) \in V},
\]
which forgets the nesting of a tuple of tuples and reindexes the entries by
the atoms of $V$; and define
\[
\psi_{F, V} : (R_Q F)(V) \longrightarrow (R_A R_R F)(V),
\qquad
\psi_{F, V}\bigl( (y_\omega)_{\omega \in V} \bigr)
= \Bigl( \bigl( y_{(i, j)} \bigr)_{j : (i, j) \in V}
  \Bigr)_{i : V \cap A_i \neq \varnothing},
\]
which groups the atomic coordinates first by ACBE block and then by RPM
block. Since the index sets on the two sides are in the canonical bijection
just described and both maps merely transport entries along that bijection,
$\varphi_{F, V}$ and $\psi_{F, V}$ are mutually inverse. At the empty context
both sides are the empty product $\mathbf{1}$ ($I_A(\varnothing) =
\varnothing$ on the nested side; no atoms on the flat side), and
$\varphi_{F, \varnothing}$ is the unique map between one-element sets: the
null cover requires no exception to the theorem.

\emph{Naturality in $V$.} Let $U \subseteq V$. On the flat side, the
restriction map of $R_Q F$ drops the coordinates of atoms
$\omega \in V \setminus U$ and applies the $F$-restriction
$F(\{\omega\}) \to F(\{\omega\})$ along the identity inclusion
$\{\omega\} \subseteq \{\omega\}$ --- which is the identity map, by
functoriality of $F$ --- on the retained atoms. On the nested side, the outer
restriction drops the $i$-coordinates with $U \cap A_i = \varnothing$ and
applies the restriction of $R_R F$ from $V \cap A_i$ to $U \cap A_i$ on the
rest; that inner restriction in turn drops the $j$-coordinates with
$U \cap A_i \cap R_j = \varnothing$ and applies the identity
$F(\{(i,j)\}) \to F(\{(i,j)\})$ on the rest. Both routes therefore implement
the same operation on entries: delete the entries indexed by atoms outside
$U$, keep the remaining entries unchanged. Hence
$\varphi_{F, U} \circ \mathrm{res}^{R_A R_R F}_{V, U}
= \mathrm{res}^{R_Q F}_{V, U} \circ \varphi_{F, V}$: flattening commutes with
restriction, and $\varphi_F$ is a map of presheaves (with inverse
$\psi_F$).

\emph{Naturality in $F$.} For a presheaf map $\tau : F \to F'$, the maps
$(R_A R_R \tau)_V$ and $(R_Q \tau)_V$ both act by applying
$\tau_{\{\omega\}}$ independently to each atomic coordinate --- on the nested
side after unwinding the two componentwise definitions, on the flat side
directly. Transporting entries along the index bijection commutes with
applying the same map to each entry, so
$\varphi_{F', V} \circ (R_A R_R \tau)_V = (R_Q \tau)_V \circ \varphi_{F, V}$:
the isomorphism is natural in $F$.

\emph{Unit compatibility.} The unit of the composite reflector at $F$ is the
composite $\eta_{A, R_R F} \circ \eta_{R, F} : F \to R_R F \to R_A R_R F$.
Chasing $x \in F(V)$: first
$\eta_{R, F, V}(x) = (x|_{V \cap R_j})_j$; then applying $\eta_{A}$ at
$R_R F$ restricts this tuple to each ACBE piece, and the restriction of
$R_R F$ from $V$ to $V \cap A_i$ sends $(x|_{V \cap R_j})_j$ to
$\bigl( x|_{V \cap A_i \cap R_j} \bigr)_{j : V \cap A_i \cap R_j \neq
\varnothing}$. The resulting nested tuple has $(i, j)$-entry
$x|_{\{(i, j)\}}$, so
\[
\varphi_{F, V}\bigl( (\eta_{A, R_R F} \circ \eta_{R, F})_V (x) \bigr)
= \bigl( x|_{\{\omega\}} \bigr)_{\omega \in V}
= \eta_{Q, F, V}(x):
\]
the flattening carries the composite unit to the unit of $R_Q$ exactly.

\emph{Symmetry, commutation, and identification with the join.} Exchanging
the roles of $A$ and $R$ throughout gives a unit-compatible natural
isomorphism $R_R R_A \cong R_Q$ by the same three checks (the objectwise
computation is symmetric in the two partitions because transversality
\eqref{ch10:eq:transverse} is). Composing the first comparison with the
inverse of the second yields a canonical natural isomorphism
$R_A R_R \cong R_R R_A$. By Corollary~\ref{ch10:cor:thereflector},
$M_A M_R \cong R_A R_R$ and $M_R M_A \cong R_R R_A$ and
$M_{\mathrm{atom}} \cong R_Q$; and by Lemmas~\ref{ch10:lem:universal}
and~\ref{ch10:lem:join}, $R_Q$ is the associated-sheaf reflector for
$J_{\mathrm{atom}} = J_A \vee J_R$, so either composite is naturally
isomorphic to sheafification for the joined topology.

Finally, the isomorphisms are canonical --- they are determined by the index
bijection induced by transversality, with no choices --- but they are not
strict equalities: nested and flat $\Set$-products need not be literally
equal sets in a chosen implementation, and the theorem asserts nothing
stronger than natural isomorphism.
\end{proof}

\section{Corollaries}
\label{ch10:sec:corollaries}

\begin{corollary}[Commutation]
\label{ch10:cor:commute}
The finite-context modalities commute up to canonical natural isomorphism:
$M_A M_R \cong M_R M_A$, compatibly with the composite units.
\end{corollary}

\begin{proof}
Immediate from Theorem~\ref{ch10:thm:composition}, composing the two
comparisons through $M_{\mathrm{atom}}$; unit compatibility holds because
both comparisons carry their composite units to the unit of
$M_{\mathrm{atom}}$.
\end{proof}

\begin{corollary}[Composites land in the intersection subcategory]
\label{ch10:cor:intersection}
For every presheaf $F$, the composite value $M_A M_R F$ (equivalently
$M_R M_A F$, equivalently $M_{\mathrm{atom}} F$) is both a $J_A$-sheaf and a
$J_R$-sheaf: either composite lands in
$\Sh(L, J_A) \cap \Sh(L, J_R) \subseteq \PSh(L, \Set)$.
\end{corollary}

\begin{proof}
By Theorem~\ref{ch10:thm:composition} the composite is naturally isomorphic
to $R_Q F$, a sheaf for $J_{\mathrm{atom}} = J_A \vee J_R$ by
Lemma~\ref{ch10:lem:sheafness}. Since $J_A \le J_A \vee J_R$ and
$J_R \le J_A \vee J_R$, Lemma~\ref{ch10:lem:coarser} makes every
$(J_A \vee J_R)$-sheaf a sheaf for each of the coarser topologies.
\end{proof}

\begin{corollary}[Identity on atomic sheaves]
\label{ch10:cor:atomicfix}
On an already $J_{\mathrm{atom}}$-sheaf, both modalities and both composites
are naturally isomorphic to the identity. This does \emph{not} make either
modality the identity on the ambient presheaf category: the witnesses $F_A$
and $F_R$ are moved (indeed, by Example~\ref{ch10:ex:action}(2), each is
collapsed to the terminal presheaf by the opposite reflector).
\end{corollary}

\begin{proof}
Let $F$ be a $J_{\mathrm{atom}}$-sheaf. By Lemma~\ref{ch10:lem:coarser}, $F$
is also a $J_A$-sheaf and a $J_R$-sheaf, so by
Corollary~\ref{ch10:cor:fixed} each unit $\eta_{A, F}$, $\eta_{R, F}$,
$\eta_{Q, F}$ is an isomorphism; the composite units are then isomorphisms as
well, and unit compatibility (Theorem~\ref{ch10:thm:composition}) transports
this to either composite. The non-triviality clause is the distinctness
theorem (Theorem NT2-3, proved earlier in the monograph) together with
Example~\ref{ch10:ex:action}(2).
\end{proof}

\begin{theorem}[Identity on the resolved subcategory]
\label{ch10:thm:resolvedidentity}
Every sheaf for $J_A \vee J_R$ is separately a $J_A$-sheaf and a
$J_R$-sheaf. Consequently both reflectors $M_A$ and $M_R$ restrict, up to
their canonical reflector isomorphisms, to the identity on the resolved
subcategory $\Sh(L, J_A \vee J_R)$; and the same statement transports along
the ambient design of Chapter~12 to the principle that no text may describe
$M_A$ or $M_R$ as nontrivial endomodalities of the resolved subtopos
$\Eres$.
\end{theorem}

\begin{proof}
The first sentence is Lemma~\ref{ch10:lem:coarser} applied to
$J_A \le J_A \vee J_R$ and $J_R \le J_A \vee J_R$: a sheaf for the finer
topology satisfies the gluing bijection for every covering sieve of any
coarser topology contained in it. For the second sentence: a reflector is
naturally isomorphic to the identity exactly on its subcategory of local
objects (Corollary~\ref{ch10:cor:fixed}), and the first sentence places every
resolved object in both subcategories of local objects simultaneously. The
final clause is a design consequence, not an additional theorem: the
nontrivial modal action of $M_A$ and $M_R$ exists only on objects not yet
resolved, which is precisely why the successor package of Chapter~12 keeps
the ambient presheaf topos, and not the resolved subtopos, as the carrier of
the modalities.
\end{proof}

\begin{proposition}[The modality algebra]
\label{ch10:prop:algebra}
Up to canonical natural isomorphism, the set
$\{\id, M_A, M_R, M_{\mathrm{atom}}\}$ is closed under composition, and
composition realizes the join of the corresponding topologies
$\{J_{\mathrm{triv}}, J_A, J_R, J_{\mathrm{atom}}\}$ (with
$J_{\mathrm{triv}}$ the minimal topology, whose modality is the
identity): for topologies $J, J'$ in this set with modalities $M, M'$,
\[
M M' \;\cong\; M_{J \vee J'} .
\]
In particular $M_A M_A \cong M_A$, $M_R M_R \cong M_R$,
$M_A M_{\mathrm{atom}} \cong M_{\mathrm{atom}} M_A \cong
M_{\mathrm{atom}}$, and symmetrically for $M_R$; the only composite
producing something new from distinct nontrivial inputs is
$M_A M_R \cong M_R M_A \cong M_{\mathrm{atom}}$.
\end{proposition}

\begin{proof}
The mixed case $M_A M_R$ is Theorem~\ref{ch10:thm:composition}, and
composites with the identity are trivial. Idempotency
$M_P M_P \cong M_P$ is Corollary~\ref{ch10:cor:fixed}. For
$M_{\mathrm{atom}} M_A \cong M_{\mathrm{atom}}$, compute with the
explicit formulas:
\[
(R_Q R_A F)(V) = \prod_{\omega \in V} (R_A F)(\{\omega\})
= \prod_{\omega \in V} F(\{\omega\}) = (R_Q F)(V),
\]
since a singleton context meets exactly one ACBE block and the
singleton-factor bijection applies; the identification is natural in
$V$ and $F$ and unit-compatible by the same coordinate bookkeeping as
in Theorem~\ref{ch10:thm:composition}. For
$M_A M_{\mathrm{atom}} \cong M_{\mathrm{atom}}$:
\[
(R_A R_Q F)(V) = \prod_{i \,:\, V \cap A_i \neq \varnothing}
(R_Q F)(V \cap A_i)
= \prod_{i} \prod_{\omega \in V \cap A_i} F(\{\omega\})
\;\cong\; \prod_{\omega \in V} F(\{\omega\}),
\]
regrouping the atoms of $V$ by the ACBE block containing each --- a
special case of the flattening argument. The RPM cases are mirrors.
Finally, the composition table matches the join table of the four
topologies: $J_A \vee J_A = J_A$, $J_A \vee J_{\mathrm{atom}} =
J_{\mathrm{atom}}$ (by $J_A \le J_{\mathrm{atom}}$, Lemma
\ref{ch10:lem:join} Step 1), and $J_A \vee J_R = J_{\mathrm{atom}}$
(Lemma~\ref{ch10:lem:join}).
\end{proof}

\begin{remark}
\label{ch10:rem:algebra}
Proposition~\ref{ch10:prop:algebra} packages the chapter's results as a
single algebraic statement: the finite-context modalities form a
commutative idempotent monoid canonically isomorphic to the
join-semilattice generated by $J_A$ and $J_R$. That semilattice has
exactly four elements --- the trivial topology, the two factor
topologies, and their join --- so the modal landscape of the finite
site is completely mapped: there are no hidden composites, no
order-dependent operators, and no further modalities generable from
these by composition. Any richer modal algebra requires a richer site.
\end{remark}

\begin{example}[Both routes, on concrete fixtures]
\label{ch10:ex:bothroutes}
Two spot computations, both reproduced by the executable layer of
Chapter~11, illustrate the composition theorem numerically.

\emph{On the ACBE witness.} At the global context,
$(R_R F_A)(\Omega) = F_A(R_0) \times F_A(R_1) = \mathbf{1} \times
\mathbf{1}$, and more generally $R_R F_A$ is terminal
(Example~\ref{ch10:ex:action}(2), with roles exchanged); applying
$R_A$ to a terminal presheaf returns a terminal presheaf. Along the
other route, $(R_Q F_A)(V) = \prod_{\omega \in V} F_A(\{\omega\})$ is
a product of one-element sets. Both composites are terminal, as the
theorem requires --- and the example shows the composite destroying
strictly more than either factor preserves: $F_A$ is fixed by $M_A$
and annihilated by $M_R$, so the composite in either order
annihilates it.

\emph{On the constant two-section presheaf.} Let $C$ be constant at
$\{0, 1\}$. Then $(R_R C)(A_i) = C(\{(i,0)\}) \times C(\{(i,1)\}) $
has $4$ elements, so $(R_A R_R C)(\Omega) = 4 \times 4 = 16$; and
directly $(R_Q C)(\Omega) = \prod_{\omega \in \Omega} C(\{\omega\})
= 2^4 = 16$. At a triple $T_{ij}$ both routes give $2^3 = 8$; at the
empty context both give $1$ (against $C(\varnothing) = \{0,1\}$,
exhibiting the empty-context correction). The flattening map matches
the nested $4 \times 4$ indexing with the flat $2^4$ indexing atom by
atom, which is exactly the bijection whose naturality the verifier
checks on this fixture.
\end{example}

\section{The narrative correction: distinct but commuting}
\label{ch10:sec:narrative}

\begin{finding}
\label{ch10:fnd:narrative}
The program's earlier informal materials hoped that the ACBE and RPM lenses
would eventually support \emph{interacting, order-sensitive} modal structure
--- that resolving evidence through the ACBE lens first and the RPM lens
second might differ from the opposite order, and that this order sensitivity
might itself carry diagnostic content. Theorem~\ref{ch10:thm:composition}
closes that hope for the present finite site, in both directions at once. The
modalities are \emph{distinct} and \emph{non-identity}: Theorem NT2-3
exhibits explicit witnesses separated by them, and
Example~\ref{ch10:ex:action}(2) shows each can annihilate the other's
witness. Yet they \emph{commute}: both composites are canonically the joined
sheafification, and the order of application is invisible up to canonical
natural isomorphism.
\end{finding}

\begin{boundary}
\label{ch10:bnd:noncommutation}
Distinctness does not imply noncommutation, and no current artifact licenses
a noncommutation claim. Any text, in any register of the program, that
describes the ACBE and RPM modalities of the present finite site as
noncommuting, order-sensitive, or ``interacting'' misstates the registered
mathematics. A successor construction seeking genuinely order-sensitive
modalities must \emph{change the site again} --- different generators, a
different context object, or added structure --- rather than reinterpret this
one; the present chapter proves that on this site the order sensitivity is
not there to be found. Symmetrically, the commutation proved here is a
finite-context fact about $\PSh(L, \Set)$ only: the product-site lift to
$\PSh(\cC \times L)$, its coherence in the $\cC$ coordinate, and its
interaction with the observational topology remain \textsc{open}
(Chapter~12), so no artifact may assert ambient commutation either.
\end{boundary}

\begin{remark}[What the correction changes downstream]
\label{ch10:rem:downstream}
Three downstream consequences are worth recording explicitly. First, any
pipeline design that schedules ACBE resolution before RPM resolution (or vice
versa) for \emph{mathematical} reasons is unjustified at the finite-context
level: the results agree canonically, and a scheduling choice can be defended
only on implementation or audit grounds --- for example, because the marked
provenance records of the two orders differ even though the sheaf-theoretic
outputs do not. Second, the joined resolution is maximally destructive among
the three modalities: by Corollary~\ref{ch10:cor:intersection} it factors
through both sheaf conditions, and on the witnesses it collapses everything
(both $R_Q F_A$ and $R_Q F_R$ are terminal, since no complete block survives
inside a singleton). Whatever information a workflow wishes to retain must be
retained \emph{before} joint resolution or outside the topology, in the
marked records. Third, the resolved subcategory is modally inert
(Theorem~\ref{ch10:thm:resolvedidentity}); the ambient category is the only
home for the modal dynamics, which is the formal content of the package
definition $\mathrm{NoeTop}_2 := \Eamb$ discussed in Chapter~12.
\end{remark}

\section{Status recapitulation}
\label{ch10:sec:status}

\begin{boundary}
\label{ch10:bnd:recap}
The results of this chapter --- Lemmas~\ref{ch10:lem:sheafness},
\ref{ch10:lem:universal}, \ref{ch10:lem:join} and
Theorem~\ref{ch10:thm:composition} with its corollaries --- constitute the
registered finite-context modality-composition record. Its status is
\textsc{internal proof candidate / review required}: an informal proof is
present in full (this chapter), exhaustive executable corroboration exists
(Chapter~11), and nothing else. The record is Ring~4 and \textsc{not
licensed}; it has not been independently reviewed, not formalized beyond the
bounded Lean scope stated in Chapter~11, and carries
\textsc{novelty unassessed} status --- no literature search has established
whether the explicit product-over-blocks reflector for transverse partition
coverages on a powerset lattice, elementary as it is, is or is not standard.
The chapter licenses no claim about $\PSh(\cC \times L)$, no probabilistic or
enriched gluing statement, no marked-object semantics (the proof concerns
ordinary $\Set$-valued presheaves; external provenance records neither enter
nor are affected), and no empirical, causal, safety, or product claim of any
kind.
\end{boundary}
